% =====================================================================
% Chapter 1 -- Introduction
%
% Redrafted from 2026-10-23 (outline: outlines/ch1_introduction.md, revision 2).
% The author's old notes are kept, commented out, at the end of this file
% until all sections are redrafted (backup: backups/introduction_2026-10-23_before_redraft.tex).
% =====================================================================

\chapter{Introduction}
\label{chap:intro}

% ---------------------------------------------------------------------
\section{Motivation}
\label{sec:intro-motivation}

% CLAUDE-DRAFT
Medical image segmentation assigns every voxel of an image either to a structure, such as an organ or a tissue, or to the background.
The resulting masks make it possible to measure a structure instead of judging it by eye.
Drawing such masks by hand is laborious, however, and different readers draw them differently~\cite{desai2021skmtea}.
In a three-dimensional (3D) volume, the structure has to be outlined on every slice in which it appears.

% CLAUDE-DRAFT
Automatic segmentation with deep learning has commonly meant one model per dataset and set of target structures, built for example with nnU-Net~\cite{isensee2018nnunetselfadaptingframeworkunetbased}.
The developers of SAT, for example, compared their single model with 72 nnU-Nets, one trained on each of their datasets~\cite{zhao2025sat}.
Universal segmentation models replace such collections with one network that segments many structures and receives, as an additional input, a prompt that specifies the target~\cite{zhao2025sat,rokuss2025voxtellfreetextpromptableuniversal}.
The prompt can be spatial, such as a point or a box marked in the image~\cite{Ma_2024,cheng2023sammed2d}, or a text that names the target (\cref{sec:text-promptable-seg}).

% CLAUDE-DRAFT
Text-promptable models differ in the text they accept.
SAT segments \gls{ct}, \gls{mri} and \gls{pet} volumes and is prompted with terms from a vocabulary of 497 classes; free text has to be mapped to these terms by a separate language model first~\cite{zhao2025sat}.
VoxTell, the model studied in this thesis, accepts a single anatomical term as well as a full clinical sentence~\cite{rokuss2025voxtellfreetextpromptableuniversal}.
We call models of this kind \emph{free-text promptable}.
In VoxTell, the structure to segment is not tied to an output channel of the network: the same weights serve every structure, and the text prompt alone selects the target~\cite[Sec.~3]{rokuss2025voxtellfreetextpromptableuniversal} (\cref{subsec:bg-voxtell}).

% CLAUDE-DRAFT
Trained on more than 62{,}000 \gls{ct}, \gls{mri} and \gls{pet} volumes, VoxTell segments structures of datasets it was not trained on with a mean \gls{dsc}, the overlap between predicted and reference mask, of 70.85\,\%~\cite[Table~1]{rokuss2025voxtellfreetextpromptableuniversal}.
On knee \gls{mri} from the SKM-TEA dataset~\cite{desai2021skmtea}, its mean \gls{dsc} over six structures, the femoral, patellar, medial tibial and lateral tibial cartilage and the medial and lateral meniscus, was 10.72\,\%~\cite[Supplementary Table~9]{rokuss2025voxtellfreetextpromptableuniversal}.
Its authors explain the result by the small share of the knee in the training data and by structures that the training data lack, and they suggest adapting the model with a few labelled examples as future work~\cite[Sec.~6]{rokuss2025voxtellfreetextpromptableuniversal}.
According to the label lists of the paper, the femoral, patellar and tibial cartilages are part of VoxTell's training data and the menisci are not~\cite[Supplementary Tables~4, 5 and~10]{rokuss2025voxtellfreetextpromptableuniversal}.
We therefore regard the knee structures as rarely seen rather than unseen, although the authors describe them as absent from training~\cite[Sec.~6]{rokuss2025voxtellfreetextpromptableuniversal}.
On the test volumes of our own data split, the pretrained model reached a mean \gls{dsc} of 0.114 over the same six structures (\cref{sec:exp-main}).

% CLAUDE-DRAFT
These structures are worth segmenting automatically.
SKM-TEA was acquired with the \gls{qdess} sequence, whose images yield maps of the T2 relaxation time, a quantity that reflects the degeneration of collagen-rich tissue, as well as quantitative biomarkers for osteoarthritis~\cite{desai2021skmtea}.
Such values are read out within the cartilage and menisci, and Desai et al.\ name the manual localisation of these regions as one of the obstacles to the use of quantitative \gls{mri}~\cite{desai2021skmtea}.

% CLAUDE-DRAFT
A model such as VoxTell would therefore have to learn these structures from a few annotated volumes, because every annotation has to be drawn by hand.
The annotations also need not arrive at once.
A group that needs several new structures may annotate them one after another and add each structure to the model as soon as its annotations are available, without losing what the model could segment before.

% ---------------------------------------------------------------------
\section{Problem Statement}
\label{sec:intro-problem}

% CLAUDE-DRAFT
We call the learning of one new target structure, named by one text prompt, a \emph{task}.
In the setting of this thesis, tasks arrive one after another, and after each task the model has to segment every structure learned so far as well as the structures it could segment before the first task.
The experiments use six knee structures as tasks (\cref{sec:intro-approach}).

% CLAUDE-DRAFT
The model therefore has to learn the new structure, keep the structures learned earlier in the sequence, and keep the ability it had before the first task, which we call its pretrained ability (\cref{fig:ch1_setting}).
The simplest way to learn a task is to fine-tune the model on its data, but training on later tasks can lower the performance on earlier ones, an effect known as catastrophic forgetting~\cite[Sec.~1]{delange2022}.
For pretrained vision--language models, continual fine-tuning also lowers the zero-shot accuracy on data the model has not been trained on~\cite{zheng2023preventingzeroshottransferdegradation}.
A strategy that changes the model little keeps what the model knew but may learn the new structure poorly, and a strategy that adapts freely risks the opposite; this conflict between stability and plasticity is a known dilemma of continual learning~\cite[Sec.~1]{delange2022}.
In this thesis, forgetting covers both losses: the loss on earlier tasks of the sequence and the loss of pretrained ability.

\begin{figure}[tb]
\centering
\includegraphics{figures/ch1_setting.pdf}
% CLAUDE-DRAFT
\caption[Setting of this thesis]{Setting of this thesis, shown for the original task order.
The pretrained model VoxTell is adapted to six structures of the SKM-TEA knee dataset one after another, each named by a text prompt and given with $k$ labelled volumes ($k = 5$, 10 or 20, or the full training split of 108 volumes); tasks~3 to~5 (femoral cartilage, patellar cartilage, medial meniscus) are not shown.
After each task, the same model has to segment the structure just learned, the structures learned before and what it could segment before, represented by the lung in AeroPath, the control task.
The reverse task order is run as well.
Images: slices of the example volumes of \cref{fig:ch4_examples} with their reference masks.}
\label{fig:ch1_setting}
\end{figure}

% CLAUDE-DRAFT
Each task comes with only a few labelled volumes.
By \emph{few-shot} we mean fine-tuning the model on $k$ labelled volumes per task, where one shot is one labelled training volume; we do not mean segmenting a new structure from a few examples given at inference without changing the weights (\cref{sec:sota-few-shot}).
Few labelled volumes are the realistic case, because each volume has to be annotated by hand (\cref{sec:intro-motivation}).
They also mean fewer updates: with the training schedule of our experiments, a task with five volumes is trained with 50 updates instead of 1{,}080 with the full training split (\cref{subsec:exp-data}).
A strategy then has less opportunity to learn the new structure, and how this shifts the balance between learning and forgetting cannot be read from results obtained with full training sets.

% CLAUDE-DRAFT
Existing continual-learning results do not show how the strategies behave in this setting.
The continual-learning methods for pretrained vision--language models reviewed in \cref{sec:cl-vlm} were developed and evaluated on classification or on tasks with text output.
Continual learning for medical image segmentation has mostly trained networks for the tasks or given each new task components of its own, such as a decoder or an output head~\cite{ji2023continualsegmentsingleunified,guo2025continuallearningdrivenmodelaccurate,zhang2023continuallearningabdominalmultiorgan,ye2025medseqft} (\cref{sec:cl-medseg}).
A free-text promptable model needs no new output for a new structure, because the text prompt names the structure and the same weights produce its mask (\cref{subsec:bg-voxtell}).
Whether a strategy should then change the shared weights, add parameters of its own, or constrain the update by replay or distillation is not settled by these results.

% CLAUDE-DRAFT
In our search of the literature (\cref{sec:sota-summary}), we found no work that adapts a pretrained free-text promptable 3D segmentation model to a sequence of new structures with few labelled volumes per task.
The closest works adapt the 2D model CLIPSeg with full training sets~\cite{gao2026medcrpcl,gao2025paewc,gao2026medpeftcl} or adapt a 3D model with few samples but without text prompts~\cite{chen2026fewshotcontinuallearning3d}.
None of them measures how much of the pretrained model's ability outside the task sequence is kept after the last task.

% ---------------------------------------------------------------------
\section{Research Question and Objectives}
\label{sec:intro-rq}

% CLAUDE-DRAFT
This thesis addresses the following research question:
\begin{quote}
\textbf{How do different continual-learning strategies trade off adaptation performance and catastrophic forgetting when adapting free-text promptable universal 3D medical image segmentation models under few-shot conditions?}
\end{quote}
Forgetting is meant in the sense of \cref{sec:intro-problem}: the loss on earlier tasks of the sequence and the loss of pretrained ability.
Few-shot conditions are $k = 5$, 10 and 20 labelled volumes per task, which we compare with the full training split; \cref{sec:intro-approach} states which strategies were run with few volumes.

% CLAUDE-DRAFT
We pursue four objectives:
\begin{enumerate}
  \item Develop a few-shot continual-learning evaluation protocol for free-text promptable universal 3D medical image segmentation (\cref{sec:method-protocol}).
  \item Compare continual-learning strategies under this protocol in terms of their adaptation and forgetting (\cref{sec:exp-main}).
  \item Investigate how the amount of labelled training data influences the evaluated strategies, separating the effect of fewer volumes from that of fewer updates (\cref{sec:exp-few-shot}).\todo{rerun: objective 3 relies on the equal-updates runs ($k = 5$, 1{,}080 updates per task); reword if they are not available by the cutoff}
  \item Analyse how the task order, the LoRA rank and the components of the strategies influence adaptation and forgetting (\cref{sec:exp-rank,sec:exp-task-order,sec:exp-components}).\todo{rerun: components = replay; add weight ensemble and gradient scaling if their reruns are used}
\end{enumerate}

% ---------------------------------------------------------------------
\section{Approach and Scope}
\label{sec:intro-approach}

% CLAUDE-DRAFT
We study the research question with VoxTell~\cite{rokuss2025voxtellfreetextpromptableuniversal} as the pretrained model and six structures of the SKM-TEA knee \gls{mri} dataset~\cite{desai2021skmtea} as tasks.
The tasks are trained in two task orders, an original order chosen at random and its reverse, to see whether the results depend on the order (\cref{subsec:exp-data}).
The lung, segmented in the chest \gls{ct} volumes of AeroPath~\cite{støverud2023aeropathairwaysegmentationbenchmark} (\cref{sec:data-aeropath}), serves as a control task: it is not trained on and is used only to measure how much of the pretrained ability is kept, with one exception, the replay memory of one strategy, which contains control-task training volumes (\cref{sec:method-protocol}).\todo{rerun: remove the exception if the replay rerun without control-task volumes replaces the existing run}
After each task, the model is tested on all six tasks and on the control task (\cref{sec:method-protocol}).

% CLAUDE-DRAFT
Sequential fine-tuning, which trains all parameters on each task without any continual-learning mechanism, is the baseline.
Task-specific LoRA trains one low-rank adapter~\cite{hu2022lora} per task, here with rank $r = 4$, and selects it by the task identity at inference.
Shared LoRA trains one adapter of rank $r = 1024$ for all tasks, and shared LoRA with replay adds a memory of training samples from earlier tasks.
The distillation-based strategy (adapted from ZSCL~\cite{zheng2023preventingzeroshottransferdegradation}) trains all parameters and distils the features of the frozen pretrained model, and the prompt-tuning strategy (adapted from CPE-CLIP~\cite{D_Alessandro_2023}) freezes the pretrained weights and trains learned prompt tokens together with a network that shifts the encoded text prompt, our addition inspired by CoCoOp~\cite{zhou2022cocoop}.
Together, the strategies cover replay, regularisation and parameter-efficient adaptation, the three groups of continual-learning methods for vision--language models~\cite[Sec.~4]{liu2026continuallearningvlmssurvey} reviewed in \cref{sec:cl-vlm}, and all except task-specific LoRA use one parameter set for all tasks (\cref{sec:method-strategies}).

% CLAUDE-DRAFT
All strategies are trained with the full training split of 108 volumes per task.
Sequential fine-tuning, task-specific LoRA and shared LoRA are also trained with $k = 5$, 10 and 20 volumes per task, with one draw of volumes for each~$k$ (\cref{subsec:exp-data}).\todo{rerun: update the list of strategies with few-shot runs and add the equal-updates runs after the reruns (cutoff 2026-11-01); same sentence as in \cref{sec:sota-summary}}\todo{defense: the research question names few-shot conditions, but the distillation-based and prompt-tuning strategies are run only with the full training split (D46)}

% CLAUDE-DRAFT
From the test results we compute three primary measures (\cref{sec:bg-metrics}).
Adaptation is the mean \gls{dsc} on each task right after training on it.
Final Forgetting is how far the \gls{dsc} on the earlier tasks lies below its best earlier value after the last task.
The control-task drop is the decrease of the \gls{dsc} on the control task from the pretrained model to the model after the last task.
The pretrained model, evaluated in the same way, serves as a reference.

% CLAUDE-DRAFT
This thesis does not propose a new continual-learning method; it adapts existing ones to VoxTell and compares them.
The study uses one model, one dataset for the tasks, one control task, one fixed text prompt per structure and one draw of training volumes per~$k$, and each configuration was trained once, so its conclusions are limited to this setting (\todo{ref: Ch. 7, limitations}Chapter~7).
Clinical deployment and studies with clinical users are outside its scope; the training cost of the strategies is reported in \cref{sec:exp-compute}.

% ---------------------------------------------------------------------
\section{Contributions}
\label{sec:intro-contributions}
\todo{provisional: finalise the contributions after the results analysis}

% CLAUDE-DRAFT
The contributions of this thesis are:
\begin{enumerate}
  \item An evaluation protocol for few-shot continual adaptation of a free-text promptable 3D segmentation model, which measures adaptation, forgetting within the task sequence and the loss of pretrained ability on a control task outside the sequence, and a systematic evaluation of continual-learning strategies under this protocol (\cref{sec:method-protocol,sec:exp-main}).
  \item A comparative analysis of the adaptation, forgetting, transfer and retention of these strategies with the full training split and with few labelled volumes per task (\cref{sec:exp-main,sec:exp-few-shot}).
  \item An analysis of factors that affect continual adaptation in this setting: the amount of labelled data, the task order and the adaptation configuration, that is, the LoRA rank and the components of the strategies (\cref{sec:exp-few-shot,sec:exp-rank,sec:exp-task-order,sec:exp-components}).
\end{enumerate}

% CLAUDE-DRAFT
Sequential fine-tuning, LoRA, ZSCL and CPE-CLIP are existing methods.
Their application to VoxTell, including the changes described in \cref{sec:method-strategies}, the evaluation protocol and all experiments are the work of this thesis.

% ---------------------------------------------------------------------
\section{Thesis Outline}
\label{sec:intro-outline}

% CLAUDE-DRAFT
\Cref{chap:sota} reviews text-promptable medical image segmentation, few-shot adaptation of segmentation models and continual learning for pretrained vision--language models and for medical image segmentation, and states the gap that this thesis addresses.
\Cref{chap:background} introduces VoxTell, the continual-learning problem in formal terms, the mechanisms behind the strategies and the evaluation metrics.
\Cref{chap:data} describes the SKM-TEA and AeroPath data, the tasks and the data splits.
\Cref{chap:method} defines the evaluation protocol, the adaptation of VoxTell to one task and the strategies, together with their complexity.
\Cref{chap:experiments} gives the experimental setup and the results of the comparison of the strategies and of the experiments on the LoRA rank, the amount of labelled data, the task order and the strategy components, as well as the training cost.
Chapter~7\todo{ref: Ch. 7 label} discusses the results across the experiments, answers the research question and states the limitations, and Chapter~8\todo{ref: Ch. 8 label} summarises the thesis.
The appendix lists the training and test volumes (\cref{sec:app-volumes}) and gives the evaluation matrices of all runs (\cref{sec:app-matrices}).

% =====================================================================
% AUTHOR'S OLD NOTES (2026-10-22 and earlier), not part of the thesis.
% Delete once all sections above are drafted.
% =====================================================================
% \chapter{Introduction}
%
% Medical image segmentation aims to delineate anatomical structures or pathological regions and produce masks that support the analysis of anatomy, disease, and treatment-related findings. Deep learning-based segmentation has progressed from task-specific models developed for individual datasets towards text-promptable universal models designed to support segmentation across multiple structures, modalities, and tasks \cite{zhao2025sat, Zhao_2024, text3dsam, rokuss2025voxtellfreetextpromptableuniversal}. However, text-promptable universal models can still exhibit reduced performance on previously unseen tasks, particularly when the target task or imaging characteristics differ substantially from those represented during pretraining \cite{rokuss2025voxtellfreetextpromptableuniversal, zhao2025sat}. Adapting the model to such tasks therefore becomes necessary.
%
% This thesis investigates this problem from a continual-learning perspective, focusing on free-text promptable universal models for three-dimensional (3D) medical image segmentation. The central challenge is to adapt the model to newly introduced tasks while preserving its performance on previously learned tasks, particularly when only a small number of labelled volumes are available. The following sections establish the background and problem formulation for this study.
%
% \section{Universal Models for Medical Image Segmentation}
%
% Three-dimensional medical imaging is obtained using modalities such as computed tomography (CT), magnetic resonance imaging (MRI), and positron emission tomography (PET), which commonly produce volumetric data composed of spatially aligned slices. Unlike two-dimensional images, these volumes preserve spatial context across adjacent slices, which must be considered when segmenting anatomical structures or pathological regions \cite{DAYARATHNA2024103046}. 
%
% Deep learning is a widely used approach for automated medical image segmentation. Convolutional neural networks (CNNs), including U-Net-based architectures, have been extensively adapted to 3D medical imaging \cite{3d_unet, li2018hdenseunethybriddenselyconnected, wang2020nonlocalunetbiomedicalimage}. Many of these approaches are specialised for a particular combination of imaging modality, dataset, and target structure. The nnU-Net framework provides a systematic implementation of this task-specific paradigm by automatically adapting its configuration and training pipeline to a given dataset \cite{isensee2018nnunetselfadaptingframeworkunetbased}. Although task-specific models can achieve strong segmentation performance within their intended settings, deploying separate models for an increasing number of anatomical structures, modalities, and datasets becomes less practical \cite{zhao2025sat}.
%
% Universal medical image segmentation models aim to address this limitation by allowing a common model to produce different segmentation outputs for different target structures or regions specified through an input prompt. Early approaches to promptable universal segmentation models relied primarily on visual interaction mechanisms, requiring users to provide spatial prompts, such as points or bounding boxes, to specify the segmentation target \cite{Ma_2024, du2025segvoluniversalinteractivevolumetric, cheng2023sammed2d}. Text-promptable segmentation \cite{carion2026sam3segmentconcepts} offers clinicians a more intuitive alternative by allowing them to extract anatomical structures or pathological regions using simple natural language descriptions, such as 'segment the left ventricle' and obtain the corresponding mask. 
%
% However, medical imaging systems deployed in practice may encounter changes in data distributions as new imaging conditions, anatomical targets, pathologies, or acquisition settings are introduced \cite{qazi2026clmedicalimaging}. When multiple new tasks are introduced sequentially, adapting the model can also affect its performance on previously learned tasks \cite{kemker2017measuringcatastrophicforgettingneural}. This challenge motivates the continual-learning setting considered in this thesis.
%
%
% \section{Continual and Few-Shot Learning}
%
% One straightforward adaptation strategy is to fine-tune the pretrained model using data from the newly introduced task. However, sequential fine-tuning can cause catastrophic forgetting, where performance on previously learned tasks can decrease after subsequent updates \cite{french1999catastrophic}. For multimodal models, catastrophic forgetting can be particularly severe, with different modalities exhibiting different rates of forgetting and interactions between modalities further complicating continual adaptation \cite{yu2026recentadvancesmultimodalcontinual}.
%
% Continual learning (CL) addresses this problem by considering learning as a sequence of updates rather than a single training stage \cite{van_de_Ven_2025}. At each step, a model receives data associated with a newly introduced task and is expected to learn the corresponding capability while retaining knowledge acquired from earlier tasks. This requires balancing stability, the ability to preserve previously learned knowledge, and plasticity, the ability to acquire new knowledge. Several methods have been proposed to mitigate catastrophic forgetting in unimodal models, including regularisation-based \cite{Kirkpatrick_2017}, rehearsal-based \cite{icarl2017}, and dynamic architecture-based \cite{yoon2018lifelonglearningdynamicallyexpandable} approaches. However, applying continual learning to multimodal foundation models introduces additional challenges due to their large scale and architectural complexity \cite{liu2026continuallearningvlmssurvey}. Continual learning can be formulated under different settings, including task-incremental, class-incremental, and domain-incremental learning, which differ in how task identity, class information, and data distributions are treated across learning steps \cite{liu2026continuallearningvlmssurvey}.
%
% The continual adaptation problem becomes more challenging when only a limited number of labelled examples are available for each newly introduced task. In medical imaging, obtaining high-quality voxel-level annotations requires domain expertise and can involve substantial effort, particularly for three-dimensional volumes \cite{Litjens_2017}. Few-shot learning refers to the adaptation of a model using only a small number of labelled examples from the target task \cite{theekshana26fewshot}. Few-shot conditions are used to evaluate how continual-learning strategies behave under limited supervision.
%
% \section{Problem Statement and Motivation}
% Early work on continual learning in medical image segmentation has largely considered task-specific or relatively compact models \cite{karani2018lifelonglearningapproachbrain, memmel2021adversarialcontinuallearningmultidomain, ranem2022continualhippocampussegmentationtransformers, gonzalez2023lifelong}. Modern free-text promptable universal segmentation models, in contrast, combine large-scale visual and language representations with cross-modal interaction components. Their architectural complexity and multimodal nature introduce additional considerations for continual adaptation \cite{liu2026continuallearningvlmssurvey} that are not necessarily captured by approaches developed for conventional segmentation models.
%
% Although recent studies have begun to investigate continual learning for medical foundation models \cite{chen2026fewshotcontinuallearning3d, Zhang_2025, chen2025forgingdynamicmemoryretrievalguided}, a systematic comparison of continual-learning strategies for free-text promptable universal 3D medical image segmentation remains limited. In particular, the behaviour of different continual-learning strategies in this setting remains insufficiently understood. This question is especially relevant under few-shot conditions.
%
% \section{Research Question and Objectives}
%
% Based on the problem formulated in the previous section, this thesis investigates the following research question:
%
% \medskip
% \noindent\textbf{How do different continual-learning strategies trade off adaptation performance and catastrophic forgetting when adapting free-text promptable universal 3D medical image segmentation models under few-shot conditions?}
%
% \medskip
% \noindent The research question is addressed through the following objectives:
% \begin{enumerate}
%     \item \textbf{Develop a few-shot continual-learning evaluation protocol}
%     for free-text promptable universal 3D medical image segmentation.
%
%     \item \textbf{Compare different continual-learning strategies} under the
%     proposed evaluation protocol in terms of their adaptation and forgetting
%     behaviour.
%
%     \item \textbf{Investigate the influence of the amount of labelled training
%     data} on the performance of the evaluated strategies under few-shot
%     conditions.
%
%     \item \textbf{Analyse the influence of experimental factors} on adaptation and forgetting through quantitative experiments and ablation studies.
% \end{enumerate}
%
% \section{Contributions}
%
% The main contributions of this thesis are:
%
% \begin{enumerate}
%     \item \textbf{A systematic empirical evaluation of continual adaptation for free-text promptable universal 3D medical image segmentation.}
%
%     \item \textbf{A comparative analysis of adaptation, forgetting, transfer, and retention under limited supervision,} providing an empirical basis for understanding the behaviour of continual-learning approaches in this setting.
%
%     \item \textbf{An investigation of factors affecting continual adaptation performance,} including the amount of available labelled data, task ordering, and adaptation configuration.
%
% \end{enumerate}
%
% \section{Scope and Delimitations}
%
% The scope of this thesis is defined by the model, datasets, continual-learning
% protocol, adaptation strategies, and evaluation measures considered in the experiments. The study is limited to the VoxTell \cite{rokuss2025voxtellfreetextpromptableuniversal} model and the datasets \cite{desai2021skmtea, støverud2023aeropathairwaysegmentationbenchmark} included in the defined experimental protocol. 
%
% The continual-learning experiments are conducted using the task sequences and few-shot conditions specified in the experimental design. The study compares selected continual-learning strategies that employ different mechanisms for adaptation and mitigating catastrophic forgetting, including parameter-efficient adaptation, prompt-based adaptation, and regularisation. In addition, the study evaluates the combination of parameter-efficient adaptation with replay. The study does not aim to develop a new continual-learning algorithm or to establish a universally optimal strategy for medical image segmentation.
%
% Evaluation is based on segmentation performance and continual-learning measures that quantify adaptation to newly introduced tasks and forgetting of existing capabilities. Factors such as clinical deployment, inference latency, computational cost in clinical environments, and prospective validation with clinical users are outside the scope of this study. The experiments are conducted within the available computational and storage resources, which constrain the model configurations and experimental settings considered.
%
% Finally, the conclusions are limited to the model, datasets, task sequences, adaptation methods, and few-shot conditions evaluated in this work; their generalisability to other medical imaging tasks or foundation models remains to be established.
%
%
% % The remainder of this thesis is organised as follows. Chapter 2 reviews the state of the art in universal medical image segmentation, continual learning, and related adaptation approaches. Chapter 3 presents the theoretical background required to understand the methods and evaluation used in this thesis. Chapter 4 describes the datasets used in the experiments. Chapter 5 presents the experimental setup, results, and analysis of the investigated continual-adaptation strategies. Finally, Chapter 6 summarises the main findings of the thesis, discusses its limitations, and outlines directions for future work.
